% !TEX root = ../main.tex
\begin{ouverture}[radioactivite]{6}{Décroissance radioactive}{Partie 1 · Constitution et transformations de la matière}
{Diagramme (N,Z)\enspace·\enspace Radioactivités $\alpha$, $\beta$ et $\gamma$\enspace·\enspace Loi de décroissance\enspace·\enspace Datation et radioprotection}
\objectif{Déterminer, à partir d'un diagramme (N,Z), les isotopes radioactifs d'un élément.}
\objectif{Écrire l'équation d'une réaction nucléaire et identifier le type de radioactivité.}
\objectif{Établir et exploiter la loi de décroissance radioactive et une courbe de décroissance.}
\objectif{Expliquer le principe d'une datation et dater un événement.}
\objectif{Citer des applications médicales et des moyens de protection contre les rayonnements ionisants.}
\prerequis{Composition du noyau, symbole $^A_Z\mathrm{X}$, isotopes.}
\prerequis{Caractère aléatoire de la désintégration, demi-vie (enseignement scientifique).}
\prerequis{Équation différentielle $y'=ay$, fonctions exponentielle et logarithme.}
\end{ouverture}
\essentiel
\essSection{1}{Stabilité et instabilité des noyaux}
\begin{definition}[Noyau radioactif]
Un noyau $^A_Z\mathrm{X}$ contient $Z$ protons et $N=A-Z$ neutrons ; des isotopes ont le même nombre de protons
mais des nombres de neutrons différents. Un noyau \textbf{radioactif}, instable, se transforme spontanément en un
autre noyau en émettant une particule : c'est une \textbf{désintégration}. Elle est \textbf{aléatoire} (on ne peut
pas prévoir la date de désintégration d'un noyau donné), spontanée et inéluctable : elle ne dépend ni de la
température, ni de la pression, ni de l'espèce chimique dans laquelle le noyau est engagé.
\end{definition}
\begin{propriete}[Diagramme (N,Z) et vallée de stabilité]
Dans le diagramme (N,Z), les noyaux stables forment la \textbf{vallée de stabilité} : $N\approx Z$ pour les noyaux
légers ($Z<20$), $N>Z$ pour les noyaux plus lourds. Les noyaux qui ont trop de neutrons, à droite de la vallée,
sont émetteurs $\beta^-$ ; ceux qui ont trop de protons, à gauche, sont émetteurs $\beta^+$ ; les noyaux les plus
lourds sont souvent émetteurs $\alpha$.
\end{propriete}
\schema{06-diagramme-nz}
\begin{savoirfaire}[Déterminer les isotopes radioactifs d'un élément]
Les isotopes d'un élément occupent une même ligne horizontale du diagramme (même $Z$). Les isotopes qui ne
figurent pas parmi les noyaux stables sont radioactifs ; leur position par rapport à la vallée indique le type de
radioactivité probable.
\end{savoirfaire}
\essSection{2}{Équations des réactions nucléaires}
\begin{loi}[Lois de conservation]
Au cours d'une réaction nucléaire, il y a conservation du \textbf{nombre de nucléons} $A$ et conservation du
\textbf{nombre de charge} $Z$.
\end{loi}
\begin{propriete}[Types de radioactivité]
\par\smallskip
{\small\renewcommand{\arraystretch}{1.35}
\begin{tabularx}{\linewidth}{@{}lYll@{}}\toprule
\otTete{Type} & \otTete{Particule émise} & \otTete{Équation} & \otTete{Exemple}\\\midrule
$\alpha$ & noyau d'hélium \ce{^{4}_{2}He} & \ce{^{A}_{Z}X -> ^{A-4}_{Z-2}Y + ^{4}_{2}He} & \ce{^{222}_{86}Rn -> ^{218}_{84}Po + ^{4}_{2}He}\\
$\beta^-$ & électron \ce{^{0}_{-1}e} & \ce{^{A}_{Z}X -> ^{A}_{Z+1}Y + ^{0}_{-1}e} & \ce{^{14}_{6}C -> ^{14}_{7}N + ^{0}_{-1}e}\\
$\beta^+$ & positon \ce{^{0}_{1}e} & \ce{^{A}_{Z}X -> ^{A}_{Z-1}Y + ^{0}_{1}e} & \ce{^{18}_{9}F -> ^{18}_{8}O + ^{0}_{1}e}\\
$\gamma$ & photon de haute énergie & \ce{Y^* -> Y + $\gamma$} & \ce{^{99\text{m}}_{43}Tc -> ^{99}_{43}Tc + $\gamma$}\\\bottomrule
\end{tabularx}}
\par\smallskip
Le noyau formé est souvent dans un état excité ; il revient à son état fondamental en émettant un rayonnement $\gamma$.
\end{propriete}
\schema{06-deplacements}
\begin{savoirfaire}[Écrire l'équation d'une désintégration]
Écrire le noyau père et la particule émise, déterminer $A$ et $Z$ du noyau fils par les lois de conservation, puis
identifier l'élément grâce à son numéro atomique $Z$, dans la classification périodique.
\end{savoirfaire}
\begin{attention}
L'électron émis lors d'une désintégration $\beta^-$ ne provient pas du cortège électronique : il est créé dans le
noyau, où un neutron se transforme en proton. De même, un proton se transforme en neutron lors d'une désintégration $\beta^+$.
\end{attention}
\essSection{3}{Loi de décroissance radioactive}
\begin{definition}[Constante radioactive et activité]
Dans un échantillon contenant un grand nombre $N(t)$ de noyaux radioactifs identiques, le nombre de noyaux qui se
désintègrent pendant une durée $\mathrm{d}t$ est proportionnel à $N$ et à $\mathrm{d}t$ :
$\mathrm{d}N=-\lambda\,N\,\mathrm{d}t$. La \textbf{constante radioactive} $\lambda$, en \unit{s^{-1}}, caractérise le noyau.
L'\textbf{activité} $A=-\dfrac{\mathrm{d}N}{\mathrm{d}t}=\lambda\,N$ est le nombre de désintégrations par seconde ; elle
s'exprime en becquerels (Bq).
\end{definition}
\begin{theoreme}[Loi de décroissance radioactive]
\[N(t)=N_0\,\mathrm{e}^{-\lambda t},\qquad A(t)=A_0\,\mathrm{e}^{-\lambda t}\qquad\text{et}\qquad t_{1/2}=\frac{\ln2}{\lambda}.\]
Le temps de demi-vie $t_{1/2}$ est la durée au bout de laquelle la moitié des noyaux présents se sont désintégrés ;
au bout de $n$ demi-vies, il reste $N_0/2^n$ noyaux.
\end{theoreme}
\begin{demonstration}
La relation $\frac{\mathrm{d}N}{\mathrm{d}t}=-\lambda N$ est l'équation différentielle $y'=-\lambda y$, dont les solutions
sont les fonctions $t\mapsto C\,\mathrm{e}^{-\lambda t}$ ; la condition initiale impose $C=N_0$. En multipliant par
$\lambda$, on obtient $A(t)$. À la date $t_{1/2}$, $N=N_0/2$, d'où $\mathrm{e}^{-\lambda t_{1/2}}=\frac12$ et $\lambda\,t_{1/2}=\ln2$.
\end{demonstration}
\schema{06-decroissance}
\begin{savoirfaire}[Exploiter une loi ou une courbe de décroissance]
Lire $t_{1/2}$ à l'ordonnée $N_0/2$ (ou $A_0/2$) et en déduire $\lambda=\ln2/t_{1/2}$. La durée nécessaire pour que
l'activité soit divisée par un facteur $f$ vaut $t=\dfrac{\ln f}{\lambda}$. Le tracé de $\ln A$ en fonction de $t$ donne
une droite de coefficient directeur $-\lambda$.
\end{savoirfaire}
\begin{codecote}[0.58]
\begin{pythoncol}
import random
import matplotlib.pyplot as plt
N0, p = 1000, 0.10   # p par intervalle
N, K, L = N0, [0], [N0]
for k in range(1, 41):
    N = N - sum(1 for _ in range(N)
                if random.random() < p)
    K.append(k)
    L.append(N)
plt.plot(K, L, "o")
plt.plot(K, [N0 * (1 - p)**k for k in K])
plt.show()
\end{pythoncol}
\cote
\fontsize{9.5}{12.5}\selectfont
\capacite{Capacité numérique}\par\smallskip
Pendant chaque intervalle de temps, chacun des noyaux restants se désintègre avec la probabilité $p$ : le test
\code{random.random() < p} est vrai avec cette probabilité. La simulation reproduit la loi exponentielle
$N_0\,(1-p)^k$, avec des fluctuations aléatoires d'autant plus faibles, en valeur relative, que les noyaux sont nombreux.
\end{codecote}
\schema{06-python-simulation}
\essSection{4}{Datation, médecine et radioprotection}
\begin{propriete}[Principe d'une datation]
On mesure le nombre de noyaux radioactifs, ou l'activité, d'un échantillon et on le compare à sa valeur initiale,
supposée connue : \[t=\frac{1}{\lambda}\ln\frac{N_0}{N}=\frac{t_{1/2}}{\ln2}\ln\frac{A_0}{A}.\]
On choisit un noyau dont la demi-vie est du même ordre de grandeur que l'âge à mesurer : le carbone 14 pour des
restes organiques de moins de \num{50000} ans, l'uranium 238 ou le potassium 40 pour des roches.
\end{propriete}
\begin{exemple}
Un être vivant échange sans cesse du carbone avec l'atmosphère : l'activité du carbone 14 y reste voisine de
\qty{0,23}{Bq} par gramme de carbone. Après sa mort, elle décroît avec une demi-vie de \num{5700} ans. Un charbon de
bois d'activité \qty{0,058}{Bq.g^{-1}}, quatre fois plus faible, date donc de deux demi-vies, soit environ \num{11400} ans.
\end{exemple}
\begin{propriete}[Applications médicales]
En imagerie, un traceur radioactif administré au patient se fixe sur l'organe étudié : la scintigraphie utilise
le technétium 99m (émetteur $\gamma$, $t_{1/2}=\qty{6,0}{h}$), la tomographie par émission de positons (TEP) le
fluor 18 (émetteur $\beta^+$, $t_{1/2}=\qty{110}{min}$). En thérapie, l'iode 131 (émetteur $\beta^-$,
$t_{1/2}=\qty{8,0}{jours}$) détruit des cellules de la thyroïde ; la radiothérapie irradie les tumeurs.
\end{propriete}
\schema{06-penetration}
\begin{propriete}[Se protéger des rayonnements ionisants]
Les rayonnements $\alpha$, $\beta$ et $\gamma$ sont \textbf{ionisants} : ils peuvent endommager les cellules. On limite
l'exposition externe en réduisant la \textbf{durée} d'exposition, en augmentant la \textbf{distance} à la source et
en interposant des \textbf{écrans} adaptés au rayonnement ; on évite la \textbf{contamination} (inhalation,
ingestion), particulièrement dangereuse pour les émetteurs $\alpha$. En France, l'exposition moyenne d'un habitant est
d'environ 4,5~mSv par an, dont près d'un tiers est dû au radon.
\end{propriete}
\begin{rediger}
Pour une équation nucléaire, expliciter les conservations, par exemple « nombre de nucléons : $14=14+0$ ; nombre de
charge : $6=7+(-1)$ ». Pour une datation, écrire la loi utilisée, isoler $t$, puis donner le résultat avec un nombre
de chiffres significatifs cohérent avec les données.
\end{rediger}
\exercices

\rubrique{Découverte}
\begin{blocLivre}{ex-6-1}
\begin{exercice}[1]{Lire un extrait du diagramme (N,Z)}
Le document représente une partie du diagramme (N,Z) des éléments bore (B), carbone (C), azote (N) et oxygène (O).
\schema{06-ex01}
\begin{questions}
\item Combien de protons et de neutrons le noyau de carbone 14 contient-il ? Écrire son symbole $^A_Z\mathrm{X}$.
\item Citer les isotopes du carbone représentés. Lesquels sont stables ?
\item Quel est le type de radioactivité du carbone 11 ? du carbone 14 ? Relier ce type à la position de chaque noyau
par rapport aux noyaux stables.
\item Écrire l'équation de la désintégration du carbone 14, qui conduit à un noyau d'azote 14.
\end{questions}
\end{exercice}
\end{blocLivre}
\begin{blocLivre}{ex-6-2}
\begin{exercice}[1]{Écrire des équations de désintégration}
Les noyaux suivants sont radioactifs.
\begin{questions}[label=\otLabel{\Alph*.},leftmargin=1.9em,itemsep=0pt,parsep=0.15em]
\item Le cobalt 60, \ce{^{60}_{27}Co}, émetteur $\beta^-$, utilisé autrefois en radiothérapie.
\item Le fluor 18, \ce{^{18}_{9}F}, émetteur $\beta^+$, utilisé en imagerie médicale.
\item L'américium 241, \ce{^{241}_{95}Am}, émetteur $\alpha$, présent dans d'anciens détecteurs de fumée.
\item Le tritium, \ce{^{3}_{1}H}, émetteur $\beta^-$, utilisé dans des cadrans lumineux.
\end{questions}
\begin{donnees}
Numéros atomiques : He : 2 ; O : 8 ; F : 9 ; Co : 27 ; Ni : 28 ; Np : 93 ; Pu : 94 ; Am : 95.
\end{donnees}
\begin{questions}
\item Énoncer les lois de conservation utilisées pour écrire l'équation d'une réaction nucléaire.
\item Écrire l'équation de chacune des quatre désintégrations et nommer le noyau fils.
\item Pour quelle désintégration le noyau fils compte-t-il quatre nucléons de moins que le noyau père ?
\end{questions}
\end{exercice}
\end{blocLivre}
\begin{blocLivre}{ex-6-3}
\begin{exercice}[1]{Identifier un type de radioactivité}
On considère les réactions nucléaires suivantes, incomplètes.
\begin{questions}[label=\otLabel{\Alph*.},leftmargin=1.9em,itemsep=0pt,parsep=0.15em]
\item \ce{^{222}_{86}Rn -> ^{218}_{84}Po + ...}
\item \ce{^{137}_{55}Cs -> ^{137}_{56}Ba + ...}
\item \ce{^{11}_{6}C -> ... + ^{0}_{1}e}
\item \ce{^{232}_{90}Th -> ... + ^{4}_{2}He}
\item \ce{^{99\text{m}}_{43}Tc -> ^{99}_{43}Tc + ...}
\end{questions}
\begin{donnees}
Numéros atomiques : B : 5 ; C : 6 ; N : 7 ; Ra : 88 ; Th : 90 ; U : 92.
\end{donnees}
\begin{questions}
\item Compléter chaque équation en justifiant à l'aide des lois de conservation.
\item Identifier le type de radioactivité de chaque réaction.
\item Pourquoi la réaction E ne modifie-t-elle pas la nature de l'élément ?
\end{questions}
\end{exercice}
\end{blocLivre}
\begin{blocLivre}{ex-6-4}
\begin{exercice}[1]{Lire une courbe de décroissance}
Une gélule contenant de l'iode 131 sert à traiter une hyperthyroïdie. Le graphique donne son activité en fonction du temps.
\schema{06-ex04}
\begin{questions}
\item Définir l'activité d'un échantillon radioactif et donner son unité.
\item Déterminer graphiquement le temps de demi-vie de l'iode 131.
\item En déduire la constante radioactive $\lambda$, en jour$^{-1}$.
\item Quelle est l'activité de la gélule au bout de 24~jours ? Retrouver ce résultat sans utiliser la courbe.
\end{questions}
\end{exercice}
\end{blocLivre}
\begin{blocLivre}{ex-6-5}
\begin{exercice}[1]{Compter les demi-vies}
Pour une scintigraphie osseuse, on injecte à un patient, à 8~h, une dose de technétium 99m d'activité 600~MBq.
La demi-vie du technétium 99m vaut 6,0~h ; on ne tient pas compte de l'élimination du traceur par l'organisme.
\begin{questions}
\item Quelle est l'activité du technétium à 14~h ? à 20~h ? le lendemain à 8~h ?
\item Quelle fraction de l'activité injectée reste-t-il au bout de deux jours ?
\item Pourquoi utilise-t-on en imagerie médicale des noyaux de courte demi-vie ? Citer une contrainte qui en découle.
\end{questions}
\end{exercice}
\end{blocLivre}
\rubrique{Entraînement}
\begin{blocLivre}{ex-6-6}
\begin{exercice}[2]{Établir et utiliser la loi de décroissance}
Le césium 137, émetteur $\beta^-$ de demi-vie 30,0~ans, a été dispersé lors des accidents de Tchernobyl (1986) et de
Fukushima (2011).
\begin{donnees}
Numéros atomiques : Xe : 54 ; Cs : 55 ; Ba : 56.
\end{donnees}
\begin{questions}
\item Écrire l'équation de la désintégration du césium 137.
\item Montrer que $N(t)=N_0\,\mathrm{e}^{-\lambda t}$ vérifie $\frac{\mathrm{d}N}{\mathrm{d}t}=-\lambda N$. Établir la relation
$t_{1/2}=\frac{\ln2}{\lambda}$, puis calculer $\lambda$ en an$^{-1}$.
\item Quel pourcentage des noyaux de césium 137 dispersés en 1986 subsistera en 2086 ?
\item Au bout de combien de temps l'activité d'un sol contaminé par le césium 137 aura-t-elle été divisée par 100 ?
\end{questions}
\end{exercice}
\end{blocLivre}
\begin{blocLivre}{ex-6-7}
\begin{exercice}[2]{Le détecteur de fumée}
Certains détecteurs de fumée anciens contiennent une source d'américium 241 d'activité 37~kBq. Ils sont désormais
collectés à part en fin de vie.
\begin{donnees}
Demi-vie de l'américium 241 : 432,6~ans ; \mbox{1 an $=\qty{3,156e7}{s}$} ; $M=\qty{241}{g.mol^{-1}}$ ;
$N_\mathrm{A}=\qty{6,02e23}{mol^{-1}}$.
\end{donnees}
\begin{questions}
\item Calculer la constante radioactive $\lambda$ de l'américium 241, en \unit{s^{-1}}.
\item En déduire le nombre de noyaux d'américium 241 contenus dans la source, puis leur masse.
\item Quelle est l'activité de la source au bout de 10~ans ? Commenter.
\item Pourquoi ces détecteurs sont-ils sans danger en fonctionnement normal, et pourquoi doivent-ils pourtant être collectés à part ?
\end{questions}
\end{exercice}
\end{blocLivre}
\begin{blocLivre}{ex-6-8}
\begin{exercice}[2]{Dater les restes d'un foyer préhistorique}
L'activité du carbone 14 d'un charbon de bois provenant d'un foyer préhistorique vaut 0,060~Bq par gramme de carbone.
Celle d'un organisme vivant vaut 0,23~Bq par gramme de carbone. Le graphique représente le rapport $A/A_0$ de
l'activité du carbone 14 d'un échantillon mort à sa valeur initiale.
\schema{06-ex08}
\begin{questions}
\item Pourquoi l'activité du carbone 14 d'un être vivant reste-t-elle constante ? Pourquoi décroît-elle après sa mort ?
\item Déterminer graphiquement l'âge du charbon de bois.
\item Retrouver ce résultat par le calcul, avec $t_{1/2}=\num{5700}$ ans.
\item Au-delà de \num{50000} ans environ, la datation par le carbone 14 n'est plus possible. Proposer une explication.
\end{questions}
\end{exercice}
\end{blocLivre}
\begin{blocLivre}{ex-6-9}
\begin{exercice}[2]{Le radon dans les habitations}
Le radon 222, gaz radioactif émetteur $\alpha$, provient de la désintégration de l'uranium contenu dans certaines
roches, notamment granitiques. Il s'accumule dans les pièces mal ventilées. Depuis 2018, le niveau de référence
fixé en France est de 300~Bq par mètre cube d'air, en moyenne annuelle.
\begin{donnees}
Demi-vie du radon 222 : 3,82~jours ; numéros atomiques : Po : 84 ; Rn : 86 ; Ra : 88.
\end{donnees}
\begin{questions}
\item Écrire l'équation de la désintégration du radon 222.
\item Combien de noyaux de radon contient un mètre cube d'air dont l'activité vaut \qty{300}{Bq} ?
\item Si l'arrivée de radon dans une pièce cessait, au bout de combien de jours l'activité serait-elle divisée par 10 ?
\item Les particules $\alpha$ sont arrêtées par la peau. Pourquoi le radon est-il néanmoins dangereux ? Quelle mesure
simple limite l'exposition des habitants ?
\end{questions}
\end{exercice}
\end{blocLivre}
\begin{blocLivre}{ex-6-10}
\begin{exercice}[2]{Simuler le hasard des désintégrations}
Le programme suivant simule un échantillon de 500 noyaux, chacun ayant la probabilité $p=0{,}05$ de se désintégrer
pendant chaque seconde.
\begin{python}
import random
N0, p = 500, 0.05
N = N0
noyaux = [N]
while N > N0 / 2:
    N = N - sum(1 for _ in range(N) if random.random() < p)
    noyaux.append(N)
print(len(noyaux) - 1, N)
\end{python}
Trois exécutions successives affichent \code{13 243}, \code{15 246} et \code{12 244}. Avec \code{N0 = 50000}, on
obtient \code{14 24438}, \code{14 24384} et \code{14 24486}.
\begin{questions}
\item \capacite{Capacité numérique}Que simule le test \code{random.random() < p} ? Quel est le rôle de la boucle \code{while} ?
\item Que représente le premier nombre affiché ?
\item Pourquoi les exécutions donnent-elles des résultats différents ? Pourquoi les écarts sont-ils plus faibles avec 50\,000 noyaux ?
\item Le modèle prévoit $N_k=N_0\,(1-p)^k$ au bout de $k$ secondes. Montrer que la demi-vie vaut
$\ln2/[-\ln(1-p)]$, la calculer et comparer aux résultats de la simulation.
\end{questions}
\end{exercice}
\end{blocLivre}
\begin{blocLivre}{ex-6-11}
\begin{exercice}[2]{L'iode 131 et la thyroïde}
Pour traiter certaines maladies de la thyroïde, on administre au patient une gélule d'iode 131 d'activité 550~MBq.
La thyroïde fixe l'iode et les électrons émis détruisent localement ses cellules. En cas d'accident nucléaire, les
populations proches peuvent recevoir des comprimés d'iodure de potassium, qui contient de l'iode stable.
\begin{donnees}
Demi-vie de l'iode 131 : 8,02~jours ; numéros atomiques : Te : 52 ; I : 53 ; Xe : 54.
\end{donnees}
\begin{questions}
\item Écrire l'équation de la désintégration $\beta^-$ de l'iode 131.
\item Calculer l'activité de l'iode 131 au bout de 8,02~jours, puis au bout de 30~jours.
\item Pourquoi un émetteur $\beta^-$ convient-il à cette thérapie, alors qu'un émetteur $\gamma$ ne conviendrait pas ?
\item Expliquer le rôle des comprimés d'iode stable en cas d'accident nucléaire.
\end{questions}
\end{exercice}
\end{blocLivre}
\rubrique{Objectif bac}
\begin{blocLivre}{ex-6-12}
\begin{exercice}[3]{Préparer un examen de TEP}
\begin{docu}{1}{Un traceur produit chaque matin}
Le fluorodésoxyglucose (FDG) est une molécule de glucose dans laquelle un groupe hydroxyle est remplacé par un atome
de fluor 18. Les cellules cancéreuses, grandes consommatrices de glucose, le concentrent. Le fluor 18, émetteur
$\beta^+$ de demi-vie 110~min, est produit dans un cyclotron, puis acheminé vers les hôpitaux.
\end{docu}
\begin{docu}{2}{Deux photons émis en sens opposés}
Le positon émis s'annihile presque aussitôt avec un électron du milieu : leur masse est convertie en énergie,
emportée par deux photons $\gamma$ de même énergie, émis dans des directions opposées.
\end{docu}
\schema{06-ex12}
\begin{donnees}
$t_{1/2}=\qty{109,7}{min}$ ; masse de l'électron et du positon : $m_\mathrm{e}=\qty{9,11e-31}{kg}$ ;
$c=\qty{3,00e8}{m.s^{-1}}$ ; $\qty{1}{eV}=\qty{1,602e-19}{J}$ ; numéros atomiques : O : 8 ; F : 9 ; Ne : 10.
\end{donnees}
\begin{questions}
\item Écrire l'équation de la désintégration du fluor 18.
\item Un patient doit recevoir une activité de 250~MBq à 11~h. La dose quitte le cyclotron à 8~h~30. Quelle doit être
son activité au départ ?
\item En utilisant la relation $E=m\,c^2$, vérifier que chaque photon emporte une énergie voisine de 511~keV.
\item \initiative Après l'examen, le patient doit éviter tout contact prolongé avec de jeunes enfants pendant
quelques heures. Quelle fraction de l'activité injectée subsiste au bout de 24~h ? Conclure.
\end{questions}
\end{exercice}
\end{blocLivre}
\begin{blocLivre}{ex-6-13}
\begin{exercice}[3]{La plus vieille roche terrestre}
\begin{docu}{1}{Des cristaux de zircon}
Des cristaux de zircon découverts à Jack Hills, en Australie, comptent parmi les plus anciens minéraux terrestres
connus. En cristallisant, le zircon incorpore de l'uranium mais pas de plomb.
\end{docu}
\begin{docu}{2}{De l'uranium 238 au plomb 206}
L'uranium 238 se transforme en plomb 206, stable, par une suite de désintégrations $\alpha$ et $\beta^-$. Les noyaux
intermédiaires ont des demi-vies très courtes devant celle de l'uranium 238 : tout se passe comme si chaque noyau
d'uranium 238 qui se désintègre donnait directement un noyau de plomb 206.
\end{docu}
\begin{donnees}
\ce{^{238}_{92}U} et \ce{^{206}_{82}Pb} ; demi-vie de l'uranium 238 : \num{4.468e9}~ans ; rapport mesuré dans un
cristal : \mbox{$N(\ce{Pb})/N(\ce{U})=0{,}979$}.
\end{donnees}
\begin{questions}
\item \initiative Déterminer le nombre de désintégrations $\alpha$ et le nombre de désintégrations $\beta^-$ qui
conduisent de l'uranium 238 au plomb 206.
\item On note $N_\mathrm{U}(t)$ et $N_\mathrm{Pb}(t)$ les nombres de noyaux d'uranium 238 et de plomb 206 dans le
cristal, $t$ étant compté depuis sa formation. Montrer que $N_\mathrm{Pb}(t)=N_\mathrm{U}(t)\left(\mathrm{e}^{\lambda t}-1\right)$.
\item En déduire l'âge du cristal de zircon.
\item Pourquoi le carbone 14 ne permettrait-il pas de dater ce cristal ?
\end{questions}
\end{exercice}
\end{blocLivre}
\begin{blocLivre}{ex-6-14}
\begin{exercice}[3]{Se protéger en médecine nucléaire}
\begin{docu}{1}{Trois règles de radioprotection}
En médecine nucléaire, un manipulateur prépare des seringues de technétium 99m, émetteur $\gamma$ de demi-vie 6,0~h.
Il applique trois règles : réduire la durée d'exposition, s'éloigner de la source, interposer un écran.
\end{docu}
\begin{docu}{2}{Débit de dose}
À 10~cm d'une seringue non protégée contenant 740~MBq de technétium 99m, le débit de dose vaut environ
\qty{1,4}{mSv.h^{-1}}. Il est inversement proportionnel au carré de la distance à la source. La limite d'exposition
du public, hors exposition naturelle et médicale, est de 1~mSv par an.
\end{docu}
\schema{06-ex14}
\begin{questions}
\item Écrire l'équation de l'émission $\gamma$ du technétium 99m. Pourquoi ce noyau convient-il bien à l'imagerie ?
\item Quel est le débit de dose à 50~cm de la seringue ?
\item La seringue est glissée dans un protège-seringue en plomb de 2,0~mm d'épaisseur. Par quel facteur le débit de
dose est-il divisé ?
\item \initiative À 10~cm de la seringue, au bout de combien de temps une personne recevrait-elle 1~mSv sans
protection ? avec le protège-seringue ? Conclure.
\item Quelle est l'activité de la seringue 24~h plus tard ?
\end{questions}
\end{exercice}
\end{blocLivre}
\begin{blocLivre}{ex-6-15}
\begin{exercice}[3]{Une pile nucléaire pour Mars}
\begin{docu}{1}{L'énergie du rover Perseverance}
Le rover Perseverance tire son énergie d'un générateur thermoélectrique à radio-isotope. Sa source contient 4,8~kg de
dioxyde de plutonium ; la chaleur libérée par les désintégrations $\alpha$ du plutonium 238 est en partie convertie
en électricité, ce qui fournit environ 110~W électriques au début de la mission.
\source{D'après la NASA, Jet Propulsion Laboratory, dossier de presse de la mission Mars 2020.}
\end{docu}
\begin{donnees}
On admet que la source contient 3,5~kg de plutonium 238 ; demi-vie : 87,74~ans ; $M=\qty{238}{g.mol^{-1}}$ ;
énergie libérée par désintégration : 5,59~MeV ; $\qty{1}{eV}=\qty{1,602e-19}{J}$ ; $N_\mathrm{A}=\qty{6,02e23}{mol^{-1}}$ ;
\mbox{1 an $=\qty{3,156e7}{s}$} ; numéros atomiques : U : 92 ; Np : 93 ; Pu : 94.
\end{donnees}
\begin{questions}
\item Écrire l'équation de la désintégration du plutonium 238.
\item Calculer l'activité de 1,0~g de plutonium 238.
\item En déduire la puissance thermique libérée par 1,0~g de plutonium 238, puis par la source du rover.
\item \initiative Estimer le rendement de la conversion en électricité. Commenter.
\item Quelle part de la puissance thermique initiale subsiste au bout de 14~ans ? Pourquoi le plutonium 238
convient-il mieux que le polonium 210, de demi-vie 138~jours ?
\end{questions}
\end{exercice}
\end{blocLivre}
